\chapter{Visual Encoding}
Visual encoding is about how we visually abstract data to convey the most important informations to the human brain.
This is formalized by defining \emph{marks} and \emph{channels}.


\section{Marks and Channels}
Marks are basic geometric elements that depict items or links.
They can be e.g, points, lines, or areas.
In contrast, channels, aka. dimensions or variables, control the appearance of marks.
Channels can be \emph{identity channels} and \emph{magnitude channels}.
Identity channels can be e.g., position, hue, or shape.
Magnitude channels can be e.g., position, value, tilt, or size.
Marks and channels are subjugated to different principles.

\begin{description}
    \item[Expressiveness Principle] The expressiveness principle states, that the visual encoding should express all of, and only, the information in the dataset attributes.
        This principle dictates that categorical attributes should map to identity channels, and that the ordered attributes should map to magnitude channels.

    \item[Effectiveness Principle] The effectiveness principle states that the importance of the attribute should match the salience of the channel---its noticeability.
        The most important attributes should be encoded with the most effective channels.
\end{description}

\noindent
Marks and channels together need the following component to form a visualization.

\subsection{Charts}
Charts are the \emph{recipe} of a graph.
It tells how to use channels and marks to visualize data.
Different chart types offer different ways to represent the same data.
When choosing channels in a chart, one can choose integral channels to highlight similarity or separable channels to highlight differences.
It is vitally important to be careful when doing this, as using them both may make a graph very hard to read.
